%!TEX program = xelatex
% 如果你有参考文献（myref.bib），请按 xelatex -> bibtex -> xelatex*2 的顺序进行编译
\documentclass[dvipsnames, svgnames,a4paper,11pt]{article}
% Share-version redaction helpers
\newcommand{\answerblank}[1][3.8cm]{%
  \par\nopagebreak[4]\noindent\textbf{答：}\par\nopagebreak[4]\vspace*{#1}}
\newcommand{\recordblank}[1]{%
  \par\noindent\rule{0.95\linewidth}{0pt}\rule{0pt}{#1}\par}
\newcommand{\privateimage}[2]{%
  \rule{#1}{0pt}\rule{0pt}{#2}}
% ----------------------------------------------------
%   中山大学物理与天文学院本科实验报告模板
%   实验 CB1 迈克尔逊干涉实验（激光干涉）
% ----------------------------------------------------

\input{settings} % 导入模板的相关设置
\usepackage{lipsum}


%---------------------------------------------------------------------
%	正文
%---------------------------------------------------------------------

\begin{document}

\scoresTable{20}{}{30}{}{30}{}{80}{}

% \infoTable{专业}{年级}{姓名}{学号}{实验时间}{教师签名}
\infoTable{\rule{2.5cm}{0.4pt}}{\rule{2.5cm}{0.4pt}}{\rule{2.5cm}{0.4pt}}{\rule{3cm}{0.4pt}}{\rule{3cm}{0.4pt}}{\rule{3cm}{0.4pt}}

\begin{center}
	\LARGE 大学物理实验 \quad 实验CB1\ 迈克尔干涉实验报告
\end{center}

\textbf{【实验报告注意事项】}
\begin{enumerate}
	\item 实验报告由三部分组成：
	\begin{enumerate}
		\item 预习报告：课前认真研读\underline{\textbf{实验讲义}}，弄清实验原理；实验所需的仪器设备、用具及其使用，完成课前预习思考题；了解实验需要测量的物理量，并根据要求提前准备实验记录表格（可以打印）。（25分）
	    \item 实验记录：认真、客观记录实验条件、实验过程中的现象以及数据。实验记录请用珠笔或者钢笔书写并签名（\textcolor{red}{\textbf{用铅笔记录的被认为无效}}）。\textcolor{red}{\textbf{保持原始记录，包括写错删除部分，如因误记需要修改记录，必须按规范修改。}}（不得输入电脑打印）；离开前请实验教师检查记录并签名。（30分）
	    \item 分析讨论：处理实验原始数据，对数据的可靠性和合理性进行分析；按规范呈现数据和结果（图、表），包括数据、图表按顺序编号及其引用；分析物理现象（含回答实验思考题，写出问题思考过程，必要时按规范引用数据）；最后得出结论。（25分）
	\end{enumerate}
	\textbf{实验报告就是将预习报告、实验记录、和数据处理与分析合起来，加上本页封面。}（80分）
	\item 实验报告下次课交给带实验的老师，最后一次实验，结束一周后交给老师。
	\item 实验注意事项：实验中使用激光，请注意激光安全，\textbf{不可直视激光}；转动激光器时，可用遮挡物挡住激光，防止对他人造成伤害。
\end{enumerate}

\subsection{实验安全注意事项}
\begin{enumerate}
	\item 实验过程中，激光光源不要随意打开关闭；
	\item 严禁用手触光学镜头的表面；
	\item 严禁用强力和斜向力旋转测微头，这样会损坏测微头或其他部件；
	\item 不要拆卸传动机构，以免影响仪器正常使用；
	\item 实验过程中，数条纹时，避免桌面的振动。
\end{enumerate}

\clearpage

\setcounter{section}{0}
\nsection{大学物理实验-}{迈克尔干涉实验}{预习报告}
	
\subsection{实验目的}
\begin{enumerate}
	\item 了解迈克尔逊干涉仪的构造、原理和调节方法；
	\item 学习用迈克尔逊干涉仪测量单色光波长的方法；
	\item 学习用迈克尔逊干涉仪测量透明薄片折射率的方法。
\end{enumerate}

\subsection{仪器用具}
\begin{table}[htbp]
	\centering
	\renewcommand\arraystretch{1.6}
	\begin{tabular}{p{0.05\textwidth}|p{0.20\textwidth}|p{0.05\textwidth}|p{0.5\textwidth}}
	\hline
	编号& 仪器用具名称 & 数量 &  主要参数（型号，测量范围，测量精度等） \\
	\hline
	1 & 精密干涉仪 & 1 & SGM-3 型，精密测微头放大倍数40倍 \\
	2 & He-Ne 激光器 & 1 & 波长理论值 $\lambda = 632.8\,\mathrm{nm}$ \\
	3 & 透明薄片 & 1 & 玻璃或有机片，折射率理论值 $n \approx 1.51$ \\
	4 & 螺旋测微计 & 1 & 量程 $0\sim25\,\mathrm{mm}$，精度 $0.01\,\mathrm{mm}$ \\
	\hline
\end{tabular}
\end{table}

\subsection{原理概述}

\subsubsection{迈克尔逊干涉仪工作原理}

从光源 $S$ 发出的光束经扩束镜扩束后，射在分束镜 $P_1$ 的半透半反射膜上，将光束分为两部分：一部分从 $P_1$ 半反射膜处反射，射向平面镜 $M_1$；另一部分从 $P_1$ 透射，射向平面镜 $M_2$。因 $P_1$ 和全反射平面镜 $M_1$、$M_2$ 均成 $45°$ 角，所以两束光均垂直射到 $M_1$、$M_2$ 上，经反射后再在观察屏 $E$ 处相遇叠加产生干涉。

补偿镜 $P_2$ 与分束镜 $P_1$ 材质、厚度完全相同，使两束光在玻璃中经历的光程相等，从而两束光的光程差仅由空气中的路程差决定。

在光路中，$M_1'$ 是 $M_1$ 被 $P_1$ 半反射膜反射所形成的虚像，迈克尔逊干涉仪产生的干涉条纹与 $M_2$ 和 $M_1'$ 之间厚度为 $2d$ 的空气薄膜所产生的干涉条纹等效。

\subsubsection{测量 He-Ne 激光波长——非定域干涉（点光源圆环）}

采用激光作为点光源时，两个虚光源 $S_1'$ 和 $S_2'$ 发出的球面波在空间处处相干，产生非定域干涉。若观察屏垂直于 $S_1'$ 和 $S_2'$ 的连线，则干涉图样为一组同心圆环。

观察屏上距中心 $E$ 点距离为 $R$ 处，两束光的光程差在 $L \gg d$ 近似下为：
\begin{equation}
    \Delta L = 2d\cos\delta
\end{equation}
其中 $\delta$ 为干涉条纹的倾角。亮纹条件为：
\begin{equation}
    2d\cos\delta = k\lambda \quad (k = 0, 1, 2, \ldots)
\end{equation}

当反射镜 $M_1$ 移动距离 $d$ 时，中心处"冒出"或"消失"的圆环数目 $N$ 满足：
\begin{equation}
    2d = N\lambda
\end{equation}
由此可得激光波长：
\begin{equation}
    \lambda = \frac{2d}{N}
    \label{eq:wavelength}
\end{equation}


\subsubsection{利用旋转样品法测量透明薄片折射率}

将折射率为 $n$、厚度为 $t$ 的透明薄片置于分束镜 $P_1$ 与反射镜 $M_1$ 之间，使薄片法线与光路方向平行（即薄片与光路垂直），然后将薄片旋转角度 $\theta$。旋转过程中光在薄片中的光程发生变化，导致两臂光程差改变，干涉圆环数目随之变化。设旋转角度为 $\theta$，圆环变化数目为 $N$，薄片厚度为 $t$，折射率为 $n$，则有：
\begin{equation}
    n = \frac{t \sin^2\theta}{2t(1-\cos\theta) - N\lambda} + \frac{1}{2} \left( 1 - \cos\theta - \frac{N\lambda}{2t} \right)
    \label{eq:refraction}
\end{equation}

实验中通过测量旋转角 $\theta$、对应的圆环变化数 $N$ 及薄片厚度 $t$，代入上式即可求得折射率 $n$。
\clearpage

\subsection{实验前思考题}

\begin{question}
	什么是光的相干性？怎样才能获得相干光？
\end{question}
\answerblank[3.8cm]

\begin{question}
	什么是"相干长度"和"相干时间"？如何计算光的相干长度和相干时间？
\end{question}
\answerblank[3.8cm]

\begin{question}
	迈克尔逊干涉仪能观察到干涉条纹的条件是什么？
\end{question}
\answerblank[3.8cm]

\clearpage
\begin{table}
	\renewcommand\arraystretch{1.7}
	\centering
	\begin{tabularx}{\textwidth}{|X|X|X|X|}
	\hline
	专业：& \rule{2.5cm}{0.4pt} & 年级：& \rule{2.5cm}{0.4pt} \\
	\hline
	姓名： & \rule{2.5cm}{0.4pt} & 学号：& \rule{3cm}{0.4pt}\\
	\hline
	室温：& \rule{2.5cm}{0.4pt} & 实验地点：& \rule{2.5cm}{0.4pt} \\
	\hline
	学生签名：& \rule{2.5cm}{0.4pt} & 评分：& \\
	\hline
	实验时间：& \rule{3cm}{0.4pt} & 教师签名：& \rule{3cm}{0.4pt} \\
	\hline
	\end{tabularx}
\end{table}

\nsection{大学物理实验-}{迈克尔逊干涉实验（激光干涉）}{实验记录}

\subsection{实验内容和步骤}

\subsubsection{调节迈克尔逊干涉仪，使产生定域等倾干涉条纹，测量 He-Ne 激光波长}

\textbf{步骤：}
\begin{enumerate}
    \item 安装并打开 He-Ne 激光器（注意不要直射眼睛），但先不安装扩束镜，调节出射激光束平行于光学平台，并使激光束从分束镜 $P_1$ 的中心附近入射；
    \item 调节可调反射镜 $M_2$ 背面的三个螺钉，使得 $M_1$ 和 $M_2$ 反射的光点的最亮处在观察屏 $E$ 上重合；
    \item 装上扩束镜（以获得点光源），此时应能在观察屏上看到等倾干涉条纹（非定域同心圆环）；
    \item 记录此时 $M_1$ 处精密测微头的读数 $d_i$，然后单向旋转精密测微头移动 $M_1$，中心每"出现"（或"消失"）50个圆环时记录一次读数，共记录6组数据（适用逐差法）；
    \item 利用逐差法求出 $M_1$ 移动的距离，结合 $N=50$，由 $\lambda=2d/N$ 求得激光波长。
\end{enumerate}

\textbf{注意事项：}精密测微头读数需除以放大倍数40得到实际移动量；旋转方向保持一致，以消除回程间隙误差；开始计数前先将测微计转一圈。

\recordblank{5cm}

\recordblank{6cm}

逐差法：取 $\Delta x_j = x_{j+3} - x_j$（$j=1,2,3$），对应圆环变化数均为 $N'=\underline{\hspace{1.5cm}}$，则 $M_1$ 实际位移 $d_j = \Delta x_j / 40$，波长 $\lambda_j = 2d_j / N'$，最终取平均值。

\subsubsection{用旋转样品法测量透明薄片的折射率}

\textbf{步骤：}
\begin{enumerate}
    \item 重复步骤1，重新调节干涉仪使之产生等倾干涉圆环；
    \item 在分束镜 $P_1$ 和反射镜 $M_1$ 之间的可旋转底座上放置透明薄片，薄片与光路垂直；
    \item 单向旋转底座，每当圆环变化数累计达到 $N=10\sim15$ 个时，记录一次旋转角度 $\theta_i$，共记录3组；
    \item 用螺旋测微计测量薄片厚度 $t$ 三次取平均值，代入折射率公式求 $n$。
\end{enumerate}

\recordblank{5cm}

\recordblank{6cm}

\recordblank{6cm}


\subsection{实验过程中遇到的问题记录}
\recordblank{5cm}
\clearpage
\begin{table}
	\renewcommand\arraystretch{1.7}
	\begin{tabularx}{\textwidth}{|X|X|X|X|}
	\hline
	专业：& \rule{2.5cm}{0.4pt} & 年级：& \rule{2.5cm}{0.4pt}\\
	\hline
	姓名：& \rule{2.5cm}{0.4pt} & 学号：& \rule{3cm}{0.4pt}\\
	\hline
	日期：& \rule{3cm}{0.4pt} & 评分：& \rule{3cm}{0.4pt}\\
	\hline
	\end{tabularx}
\end{table}

\nsection{大学物理实验-}{迈克尔干涉实验}{分析与讨论}
\subsection{测量 He-Ne 激光波长}
\textbf{数据处理：}
\recordblank{3.2cm}
\textbf{结果：}与理论值 $632.8\,\mathrm{nm}$ 比较，相对误差为
\recordblank{3.2cm}
\textbf{误差分析：}主要误差来源包括：精密测微头的回程间隙、读数时的人为计数误差（漏数或多数圆环）、桌面振动导致圆环不稳定等。
\subsection{测量透明薄片折射率（旋转样品法数据处理）}
\textbf{数据处理：}
\recordblank{3.2cm}
\textbf{结果：}$n = \bar{n}$，与玻璃理论值 $n \approx 1.51$ 比较，相对误差为
\recordblank{3.2cm}
\textbf{误差分析：}误差主要来源于：旋转角度的读数误差、薄片厚度的测量误差、圆环计数误差（每次转动太快导致圆环变化计数不准）等。
\clearpage
\subsection{实验后思考题}

\begin{question}
	1. 什么是非定域干涉？什么是定域干涉？什么是等倾干涉？什么是等厚干涉？
\end{question}
\answerblank[3.8cm]

\begin{question}
	2. 如何测量透明溶液的折射率？请自行就相关实验原理进行调研，并设计具体实验方案。
\end{question}
\answerblank[3.8cm]

\begin{question}
	3. 尝试根据图推导透明薄片折射率公式
\end{question}
\answerblank[3.8cm]

\clearpage
\appendix
\appendixpage
\addappheadtotoc
\subsection{实验报告记录}
\recordblank{5.5cm}
\subsection{桌面整理照片}
\recordblank{5.5cm}
\end{document}
